\chapter{理论力学概念骨架：约束、作用量与正则结构}
\section{D'Alembert 原理与 Lagrange 方程}
约束可分为完整/非完整、定常/非定常。用独立广义坐标表示
$\vb r_\alpha=\vb r_\alpha(q_1,\ldots,q_s,t)$，同一时刻允许的虚位移为
$\delta\vb r_\alpha=\sum_i(\partial\vb r_\alpha/\partial q_i)\delta q_i$。
理想约束力不作虚功。Newton 方程点乘虚位移并求和：
\[
\boxed{\sum_\alpha(\vb F_\alpha-m_\alpha\vb a_\alpha)\cdot\delta\vb r_\alpha=0}.
\]
利用恒等式
\[
\sum_\alpha m_\alpha\vb a_\alpha\cdot\pdv{\vb r_\alpha}{q_i}
=\dv{}{t}\pdv{T}{\dot q_i}-\pdv{T}{q_i},
\]
若保守力 $Q_i=-\partial V/\partial q_i$，各独立 $\delta q_i$ 的系数为零：
\[
\boxed{\dv{}{t}\pdv{L}{\dot q_i}-\pdv{L}{q_i}=0},\qquad L=T-V.
\]
同式也由作用量 $S=\int_{t_1}^{t_2}L\dd t$ 的驻值条件得到：分部积分并令端点
$\delta q_i=0$。循环坐标不显含于 $L$，故其共轭动量守恒。时间、空间平移和旋转对称性分别对应能量、动量和角动量守恒，这是 Noether 定理的核心图像。

\section{Hamilton 形式、Poisson 括号与相空间}
定义 $p_i=\partial L/\partial\dot q_i$，在速度 Hessian 非奇异时作 Legendre 变换
\[
H(q,p,t)=\sum_i p_i\dot q_i-L.
\]
微分并代入 Lagrange 方程：
\[
\dd H=\sum_i(\dot q_i\dd p_i-\dot p_i\dd q_i)-\pdv L t\dd t,
\]
比较系数得到
\[
\boxed{\dot q_i=\pdv H{p_i},\qquad\dot p_i=-\pdv H{q_i}}.
\]
Poisson 括号
$\{f,g\}=\sum_i(\partial_{q_i}f\,\partial_{p_i}g-\partial_{p_i}f\,\partial_{q_i}g)$，故
\[
\dot f=\{f,H\}+\pdv f t.
\]
正则变换保持 Poisson 括号和 Hamilton 方程形式；量子化中
$\{q,p\}=1$ 对应 $[\hat q,\hat p]=i\hbar$。Hamilton 流在相空间的散度
$\sum_i(\partial\dot q_i/\partial q_i+\partial\dot p_i/\partial p_i)=0$，所以相体积守恒（Liouville 定理），这也是统计系综演化的基础。

\section{小振动与刚体}
在稳定平衡点附近保留二次项：
\[
T=\frac12\dot{\vb q}^{\,T}M\dot{\vb q},\qquad
V=\frac12\vb q^{T}K\vb q.
\]
方程 $M\ddot{\vb q}+K\vb q=0$，代入 $\vb q=\vb A\ee^{i\omega t}$：
\[
(K-\omega^2M)\vb A=0,\qquad\boxed{\det(K-\omega^2M)=0}.
\]
不同本征矢是正规模；用质量矩阵正交归一化后，每一模都是独立谐振子。这与晶格振动和声子量子化完全同构。

刚体角动量为 $\vb L=I\vb\omega$，$I$ 的本征方向为主轴。随体坐标中惯性张量恒定，惯性系导数与随体系导数关系给
\[
\boxed{\left(\dv{\vb L}{t}\right)_{\rm body}+\vb\omega\times\vb L=\vb N}.
\]
主轴系分量即 Euler 方程。Euler 角描述姿态，但约定多样，回答时必须先声明转动顺序。

\pitfall{最小作用量通常实际是“作用量驻值”；$H$ 只有在常见自然系统中才等于机械能；广义动量也不总是 $m\dot q$。}
\section{核心公式与面试回答}
\[
\delta S=0\Rightarrow\text{Euler--Lagrange 方程},\quad
H=p\dot q-L,\quad \dot f=\{f,H\}+\partial_t f.
\]
\interview{“Lagrange 与 Newton 的关系”：二者动力学内容等价；Lagrange 形式把理想约束吸收到广义坐标中，并使对称性更透明。“Hamilton 为何重要”：一阶正则方程、相空间几何、统计力学与量子力学的共同语言。}
